% Reference entries, M (AN-11).  Law 1: os/iword/core/edit.c iw_ed_margins
% ([L/R Mar] D0/01 at the cursor, a [HRt] first when the cursor is not at
% the left margin; every paragraph after it re-wrapped; a line holds
% (10200 - R - L) / 120 + 1 characters), face_a.c M_FMTMARL / M_FMTMARR /
% parse_inches (inches to two decimals, an optional ", at most 8, ten input
% characters), os/iword/core/block.c (copy, cut, retrieve; iw_ed_select:
% sentence, paragraph, page), face_a.c M_MOVE / M_MOVEBLK / M_MOVENB /
% M_RETMENU / M_RETRIEVE.  Shots: ../wordperfect51-decomp mint MINT12..15,
% MINT21, MINT23, MINT29, MINT31.
% GATES: test-pubs-iword test-iword-codes test-iword-wp5-roundtrip test-iword-face test-flair-iword-block test-flair-iword-format
\Gates{test-pubs-iword test-iword-codes test-iword-wp5-roundtrip test-iword-face test-flair-iword-block test-flair-iword-format}
\begin{Feature}{Margins, Left and Right}
The left and right margins are the space between the edges of the paper
and the text.  A new document has a one inch margin on each side of an 8.5
inch page, so a line holds 66 characters.  You can change the margins from
any point in the document; the change applies from there to the end of the
document, or to the next margin setting.

To change the left and right margins,
\begin{WPSteps}
\item Move the cursor to where the new margins are to begin.
\item Press \FK{Format}{Shift-F8}, then select \Ch{Line}{1}.
\item Select \Ch{Margins}{7}.  The cursor moves to the left margin setting.
\item Type the left margin in inches, for example \Typed{1.5}, and press
\K{Enter}.  To keep the left margin as it is, press \K{Enter} alone.
\item Type the right margin in inches and press \K{Enter}, or press
\K{Enter} alone to keep it.
\item Press \FK{Exit}{F7} to return to the document.
\end{WPSteps}
\Documents{wpmenu}{Shift-F8,1,7=Margins}

\BeginReplay{ref-margins}
For example, type \TypeRun{Normal margins here.} and press \PressKey{Enter};
then press \Press{Format}{Shift-F8}, select \Choose{Line}{1} and
\Choose{Margins}{7}, type \TypeName{1.5}, press \PressKey{Enter}, type
\TypeName{2}, press \PressKey{Enter} and \Press{Exit}{F7}, and type
\TypeRun{Narrow text now begins here and wraps at the new right margin of two inches.}
\WordScreen[rows=3]{ref-margins}{Text after a margin setting of 1.5 inches left and 2 inches right}
\EndReplay

\Notes
\Topic{Code}
The margins go into the document as a \Code{[L/R Mar:\textit{left},\textit{right}]}
code, for example \Code{[L/R Mar:1.5",2"]}.  On the editing screen
InitechWord asks before deleting one (see \InRef{Codes}); deleting it
returns the text after it to the margins in effect before it.

\Topic{In the Middle of a Line}
New margins take effect at the start of a line.  If the cursor is not at the
left margin when you set them, InitechWord puts a hard return in front of
the margin code, so that the text after the cursor begins a new line under
the new margins.

\Topic{Measurements}
Type a margin in inches, with up to two decimal places, and if you like an
inch mark (\Typed{1.25"}).  A margin must be less than 9 inches.  A setting
InitechWord cannot read leaves that margin as it was.

\Topic{On the Screen}
The left edge of the editing screen is always one inch from the left edge
of the paper.  With a left margin of 1.5 inches, text starts five columns in;
the status line's Pos shows the cursor's place on the paper.  In Reveal
Codes the ruler marks the margins in effect at the cursor (see \InRef{Reveal
Codes}).

\Topic{Tabs and Indents}
The tab stops are measured from the left margin, so they move with it (see
\InRef{Tab}).

\SeeAlso{Format; Indent; Justification; Reveal Codes; Tab; Word Wrap}
\end{Feature}

\begin{Feature}{Move}
Move takes text out of one place and puts it in another, copies text, or
deletes it.  It works on a block you have made, or on a sentence, a
paragraph or a page that InitechWord selects for you.

To use Move,
\begin{WPSteps}
\item Block the text first (see \InRef{Block}), then press
\FK{Move}{Ctrl-F4}.  The status line shows
\StatusLine{Move: 1 Block; 2 Tabular Column; 3 Rectangle: 0}
\par Select \Ch{Block}{1} and go on as described in \InRef{Move, Block}.
\Or
With Block off, press \FK{Move}{Ctrl-F4}.  The status line shows
\StatusLine{Move: 1 Sentence; 2 Paragraph; 3 Page; 4 Retrieve: 0}
\par Select what to move (see \InRef{Move, Sentence/Paragraph/Page}), or
select \Ch{Retrieve}{4} to put back the text you last moved or copied (see
\InRef{Move, Retrieve}).
\end{WPSteps}
\Documents{wpkey}{Ctrl-F4}
\Documents{wpkey}{Block:Ctrl-F4}
\Documents{wpmenu}{Block:Ctrl-F4,1=Block}

\Notes
\Topic{Other Selections}
Tabular Column (2) and Rectangle (3) on the first menu, and Append (4) on
the second, return you to the document as Cancel does.

\Topic{The Move Buffer}
The text you move or copy is kept, so that you can put it in the document
again with Move, Retrieve.  It stays until you move or copy other text, even
across a new document.

\SeeAlso{Block; Move, Block; Move, Retrieve; Move, Sentence/Paragraph/Page}
\end{Feature}

\begin{Feature}{Move, Block}
After you block text and choose Block from the Move menu, InitechWord asks
what to do with it:
\StatusLine{1 Move; 2 Copy; 3 Delete; 4 Append: 0}

To move or copy a block,
\begin{WPSteps}
\item Block the text and press \FK{Move}{Ctrl-F4}; select \Ch{Block}{1}.
\item Select \Ch{Move}{1} to take the block out of the document
\Or
Select \Ch{Copy}{2} to leave it where it is.
\item The status line shows \Scr{Move cursor; press Enter to retrieve.}
Move the cursor to where the text is to go.
\item Press \K{Enter}.  The text is put in front of the cursor.
\end{WPSteps}
\Documents{wpmenu}{Block:Ctrl-F4,1,1=Move}
\Documents{wpmenu}{Block:Ctrl-F4,1,2=Copy}
\Documents{wpmenu}{Block:Ctrl-F4,1,3=Delete}
To delete the block instead, select \Ch{Delete}{3} at the second step.
Move deletes it without asking (compare \InRef{Block, Delete}).

\BeginReplay{ref-moveblock}
For example, type \TypeRun{Alpha Beta Gamma.}, press \PressKey{Home},
\PressKey{Left}, press \PressKey{Right} six times to move to \textit{Beta}\Keys{Right Right Right Right Right},
press \Press{Block}{Alt-F4} and \PressKey{Right} five times to block
\textit{Beta} and the space after it\Keys{Right Right Right Right}, press
\Press{Move}{Ctrl-F4} and select \Choose{Block}{1} and \Choose{Move}{1}:
\WordScreen[rows=1]{ref-moveblock-1}{The block is out of the document; InitechWord waits for Enter}
Press \PressKey{End} and \PressKey{Enter}:
\WordScreen[rows=1]{ref-moveblock-2}{The block moved to the end of the line}
\EndReplay

\Notes
\Topic{Attributes}
A block keeps its appearance and size when it is moved or copied: if it is
bold where it starts, its bold goes with it.  The text around the place it
came from keeps its own attributes.  Where the block is put, codes that
would change nothing are left out.

\Topic{Cursor Keys}
At \Scr{Move cursor; press Enter to retrieve.} every cursor key works, and
so does Reveal Codes.  Cancel leaves the text unretrieved (see
\InRef{Cancel}).

\SeeAlso{Block; Block, Delete; Move; Move, Retrieve}
\end{Feature}

\begin{Feature}{Move, Retrieve}
Retrieve on the Move menu puts the text you last moved or copied into the
document again, at the cursor.  Use it to put the same text in several
places, or to bring back text you left unretrieved.

To retrieve the moved or copied text,
\begin{WPSteps}
\item Move the cursor to where the text is to go.
\item With Block off, press \FK{Move}{Ctrl-F4}, then select
\Ch{Retrieve}{4}.  The status line shows
\StatusLine{Retrieve: 1 Block; 2 Tabular Column; 3 Rectangle: 0}
\item Select \Ch{Block}{1}.  The text is put in front of the cursor.
\end{WPSteps}
\Documents{wpmenu}{Ctrl-F4,4=Retrieve}
\Documents{wpmenu}{Ctrl-F4,4,1=Block}

\Notes
\Topic{Nothing to Retrieve}
If you have not moved or copied any text since InitechWord started, Block
(1) returns you to the document with nothing changed.  Tabular Column (2)
and Rectangle (3) also return you to the document.

\SeeAlso{Move; Move, Block}
\end{Feature}

\begin{Feature}{Move, Sentence/Paragraph/Page}
With Block off, Move can select the sentence, the paragraph or the page the
cursor is in, so that you need not block it yourself.

To move, copy or delete a sentence, paragraph or page,
\begin{WPSteps}
\item Put the cursor anywhere in the sentence, paragraph or page.
\item Press \FK{Move}{Ctrl-F4}.
\item Select \Ch{Sentence}{1}, \Ch{Paragraph}{2} or \Ch{Page}{3}.
InitechWord highlights it and shows \Scr{1 Move; 2 Copy; 3 Delete; 4 Append: 0}.
\item Select \Ch{Move}{1} or \Ch{Copy}{2}, move the cursor to where the
text is to go, and press \K{Enter}
\Or
Select \Ch{Delete}{3} to delete it.
\end{WPSteps}
\Documents{wpmenu}{Ctrl-F4,1=Sentence}
\Documents{wpmenu}{Ctrl-F4,1,1=Move}
\Documents{wpmenu}{Ctrl-F4,1,2=Copy}
\Documents{wpmenu}{Ctrl-F4,1,3=Delete}
\Documents{wpmenu}{Ctrl-F4,2=Paragraph}
\Documents{wpmenu}{Ctrl-F4,2,1=Move}
\Documents{wpmenu}{Ctrl-F4,2,2=Copy}
\Documents{wpmenu}{Ctrl-F4,2,3=Delete}
\Documents{wpmenu}{Ctrl-F4,3=Page}
\Documents{wpmenu}{Ctrl-F4,3,1=Move}
\Documents{wpmenu}{Ctrl-F4,3,2=Copy}
\Documents{wpmenu}{Ctrl-F4,3,3=Delete}

\BeginReplay{ref-movesent}
For example, type
\TypeRun{First sentence here. Second sentence follows.  Third one.},
put the cursor in the second sentence\Keys{Home Left Right Right Right Right Right Right Right Right Right Right Right Right Right Right Right Right Right Right Right Right Right Right Right Right Right Right Right},
press \Press{Move}{Ctrl-F4} and select \Choose{Sentence}{1}:
\WordScreen[rows=1]{ref-movesent}{Move highlights the sentence the cursor is in}
\Keys{F1}\EndReplay

\Notes
\Topic{Cancel}
Press \FK{Cancel}{F1} while the text is highlighted to leave it where it is.
The highlight goes and the cursor stays at the end of the text.

\Topic{Page}
A page is the page the status line counts as Pg: its lines from the first
to the last.  A page ends after 54 lines, or at a hard page, whose code goes
with it (see \InRef{Page Break, Hard}).

\Topic{Paragraph}
A paragraph runs from after the hard return (or hard page) before it
through its own hard return.

\Topic{Sentence}
A sentence runs from the start of the sentence through the period,
question mark or exclamation mark that ends it and up to three spaces after
it.  A hard return also ends a sentence.

\SeeAlso{Block; Move; Move, Block; Move, Retrieve}
\end{Feature}
